% CN 2.19 Exercises (p.76--84)

% ======================================================================
\begin{exercise}{1}
    Let $A$ be a finite set. How many indicator functions with domain $A$ are
    there?
\end{exercise}
\begin{solbox}
\begin{solution}
    An indicator function is defined by
    \[
        \chi_A(x) = \begin{cases} 1 & \text{if } x\in A \\
            0 & \text{if } x\notin A \end{cases}
    \]

    Let $|A|=n$. The number of functions $\chi: A\to\{0,1\}$ equals
    the number of ways to independently assign each of the $n$
    elements of $A$ a value in $\{0,1\}$, which is $2^n$.

    Equivalently, each such function corresponds to exactly one
    subset $X\subseteq A$ (namely $X=\{x\in A : \chi(x)=1\}$), so the
    number of indicator functions equals $|\mathcal P(A)| = 2^n$.
\end{solution}
\end{solbox}

\begin{exercise}{2}
    Let $A$ and $B$ be any sets. Express the indicator functions of each of the
    following in terms of $\chi_A$ and $\chi_B$.
    \begin{enumerate}[label=(\alph*)]
        \item $\overline A$
        \item $A\cap B$
        \item $A\setminus B$
        \item $A\triangle B$
    \end{enumerate}
\end{exercise}
\begin{solbox}
\begin{solution}
    \textbf{(a)}
    \[
        \chi_{\overline A}(x) = 1-\chi_A(x).
    \]

    \textbf{(b)} $x\in A\cap B \iff x\in A \wedge x\in B$, so
    \[
        \chi_{A\cap B} = \chi_A\cdot\chi_B.
    \]

    \textbf{(c)} $x\in A\setminus B \iff x\in A \wedge x\notin B$, so
    \[
        \chi_{A\setminus B} = \chi_A\cdot\chi_{\overline B}
        = \chi_A\cdot(1-\chi_B) = \chi_A - \chi_A\chi_B.
    \]

    \textbf{(d)} Since $A\triangle B = (A\setminus B)\cup(B\setminus
    A)$ is a disjoint union,
    \[
        \chi_{A\triangle B} = \chi_{A\setminus B}+\chi_{B\setminus A}
        = (\chi_A-\chi_A\chi_B)+(\chi_B-\chi_A\chi_B)
        = \chi_A+\chi_B-2\chi_A\chi_B.
    \]
\end{solution}
\end{solbox}

\begin{exercise}{3}
    If $f:A\to\R$ and $g:A\to\R$, their sum $f+g:A\to\R$ is defined by
    $(f+g)(x)=f(x)+g(x)$. What is the sum of all the indicator functions of all
    the subsets of a finite set?
\end{exercise}
\begin{solbox}
\begin{proof}
    Let $A$ be a finite set with $|A|=n$. Since function addition is
    defined pointwise, for any $x_i\in A$,
    \[
        \left(\sum_{X\subseteq A}\chi_X\right)(x_i)
        = \sum_{X\subseteq A}\chi_X(x_i).
    \]

    Since $\chi_X(x_i)=1$ if $x_i\in X$ and $0$ otherwise, this sum
    counts the number of subsets $X\subseteq A$ containing $x_i$.

    \textbf{Method 1 (direct count).} A subset containing $x_i$ is
    formed by fixing $x_i$ to be included, and choosing freely,
    independently, whether each of the remaining $n-1$ elements is
    included. This gives $2^{n-1}$ such subsets.

    \textbf{Method 2 (summing over subset sizes).} For each size
    $k=1,\dots,n$, the number of size-$k$ subsets containing $x_i$ is
    obtained by fixing $x_i$ as one member and choosing the
    remaining $k-1$ members from the other $n-1$ elements:
    $\binom{n-1}{k-1}$. Summing over all sizes,
    \[
        \sum_{X\subseteq A}\chi_X(x_i)
        = \sum_{k=1}^{n}\binom{n-1}{k-1}
        = \sum_{j=0}^{n-1}\binom{n-1}{j} = 2^{n-1},
    \]
    using the binomial identity $\sum_{j=0}^{m}\binom{m}{j}=2^m$
    with $m=n-1$.

    Both methods agree: $\left(\sum_{X\subseteq A}\chi_X\right)(x_i)
    = 2^{n-1}$ for every $x_i\in A$, independent of which $x_i$ is
    chosen. So $\sum_{X\subseteq A}\chi_X$ is the constant function
    on $A$ with value $2^{n-1}$.
\end{proof}
\end{solbox}

\begin{exercise}{4}
    Draw a Venn diagram showing each of the following sets of functions and the
    relationships between them: functions, injections, surjections, bijections,
    identity functions, binary relations, ternary relations, relations.
\end{exercise}

% ======================================================================
\begin{exercise}{5}
    Functions can be viewed as sets of ordered pairs, so we can combine them
    using set operations. Suppose $f:A\to B$ and $g:C\to D$ are functions. Which
    of the following is always a function?
    \[
        f\cap g;\quad f\cup g;\quad f\setminus g;\quad g\setminus f;\quad
        f\triangle g;\quad f\times g.
    \]
    For each that is always a function, give its definition in the usual form
    (domain, codomain, rule in terms of $f$ and $g$). For each that is not
    necessarily a function, explain why (e.g.\ with examples).
\end{exercise}
\begin{solbox}
\begin{proof}
    Recall that $R\subseteq A\times B$ is a function if and only if
    $R$ satisfies both:
    \begin{itemize}
        \item \textbf{Existence}: $\{x\in A : (x,y)\in R \text{ for
            some } y\in B\} = A$;
        \item \textbf{Uniqueness (well-definedness)}: for every
            $x\in A$, the set $R(x)=\{y\in B : xRy\}$ contains
            exactly one element.
    \end{itemize}

    \medskip
    \textbf{(i) $f\cap g$.}
    \[
        (f\cap g)(x) = \{f(x)\}\cap\{g(x)\}.
    \]
    When $f(x)=g(x)$: $(f\cap g)(x) = \{f(x)\}$.
    When $f(x)\ne g(x)$: $(f\cap g)(x) = \emptyset$ --- does not
    exist. So $f\cap g$ is \emph{not necessarily} a function; it can
    fail \textbf{existence}.

    \medskip
    \textbf{(ii) $f\cup g$.} \label{5ii}
    \[
        (f\cup g)(x) = \{f(x)\}\cup\{g(x)\}.
    \]
    When $f(x)=g(x)$: $|(f\cup g)(x)|=1$.
    When $f(x)\ne g(x)$: $|(f\cup g)(x)|=2$. So $f\cup g$ is
    \emph{not necessarily} a function; it can fail
    \textbf{uniqueness}.

    \medskip
    \textbf{(iii) $f\setminus g$.}
    \[
        (f\setminus g)(x) = \{f(x)\}\setminus\{g(x)\}.
    \]
    When $f(x)=g(x)$: $|f\setminus g|=0$ --- does not exist.
    When $f(x)\ne g(x)$: $|f\setminus g|=1$. So $f\setminus g$ is
    \emph{not necessarily} a function; it can fail
    \textbf{existence}.

    \medskip
    \textbf{(iv) $g\setminus f$.}
    By the same argument as (iii) with the roles of $f,g$
    interchanged, $g\setminus f$ is \emph{not necessarily} a
    function.

    \medskip
    \textbf{(v) $f\triangle g$.}
    \[
        (f\triangle g)(x) = \{(f\setminus g)(x)\}
        \cup \{(g\setminus f)(x)\}.
    \]
    When $f(x)=g(x)$: $|f\triangle g|=0$ --- does not exist.
    When $f(x)\ne g(x)$: $|f\triangle g|=2$. So $f\triangle g$ is
    \emph{not necessarily} a function; depending on $f,g$, it can
    fail either \textbf{existence} or \textbf{uniqueness}.

    \medskip
    \textbf{(vi) $f\times g$.}
    \[
        f\times g = \{((a,b),(c,d)) : (a,b)\in f,\ (c,d)\in g\}.
    \]
    This defines a function
    \[
        f\times g : A\times C \to B\times D, \qquad
        (f\times g)(x,y) = (f(x), g(y)).
    \]
    \textbf{Existence}: for all $(x,y)\in A\times C$, since $f$ and
    $g$ are functions, $f(x)$ and $g(y)$ must exist, so
    $(f(x),g(y))$ exists.
    \textbf{Uniqueness}: since $f(x)$ and $g(y)$ are each unique,
    the pair $(f(x),g(y))$ is unique.
    So $f\times g$ is \emph{always} a function.
\end{proof}
\end{solbox}

\begin{exercise}{6}
    Suppose $f:A\to B$ and $g:C\to D$ are functions with disjoint domains,
    $A\cap C=\emptyset$. Is the disjoint union $f\sqcup g$ always a function? If
    so, give its definition in terms of $f$ and $g$. If not, why not?
\end{exercise}
\begin{solbox}
\begin{solution}
    By \hyperref[5ii]{5(ii)} $f \cup g $ is always a function.
\end{solution}
\end{solbox}

\begin{exercise}{7}
    Give an example of a function $f:\N\times\N\to\N$ which is an injection.
    Even better, give one which is a bijection.
\end{exercise}

% ======================================================================
\begin{exercise}{8}
    A \emph{bitstring} is a member of $\{0,1\}^*$.
    \begin{enumerate}[label=(\alph*)]
        \item Give a bijection from $\{0,1\}^*$ to $\N$.
    \end{enumerate}
    Let $\{0,1\}^{*+}$ denote the set of all finite nonempty tuples of
    bitstrings, e.g.\ $(011,1,10001)$, $(10,10,10,10)$, $(\epsilon,1011)$,
    $(1010010001)$.
    \begin{enumerate}[label=(\alph*),start=2]
        \item (Challenge) Give a bijection from $\{0,1\}^{*+}$ to $\{0,1\}^*$.
    \end{enumerate}
\end{exercise}
\begin{solbox}[long]
\begin{solution}[(a)]

\paragraph{Setup: grouping by length.} For each $n\ge0$, there are $2^n$
bitstrings of length $n$. We assign each bitstring a natural number
(``index'') as follows:

\begin{center}
\begin{tabular}{c c c c}
\hline
length & bitstring & number of bitstrings & index $\in \N$ \\
\hline
$0$ & $\varepsilon$ & $1$ & $0$ \\
$1$ & $0,1$ & $2=2^1$ & $1,2$ \\
$2$ & $00,01,10,11$ & $2^2$ & $3,4,5,6$ \\
$\vdots$ & $\vdots$ & $\vdots$ & $\vdots$ \\
$n$ & $\cdots$ & $2^n$ & $\cdots$ \\
\hline
\end{tabular}
\end{center}

\paragraph{Defining $f_n$.} For $\{0,1\}^n$, the value of a string is a
member of $\{0,1,\dots,2^n-1\}$. Define $f_n$ to be the value of
$\{0,1\}^n$:
\[
    f_n : \{0,1\}^n \longrightarrow \{0,1,\dots,2^n-1\}, \qquad
    f_n(s) = f_n(s_1 s_2 \cdots s_n) = \sum_{i=1}^n s_i\cdot 2^{n-i}.
\]

\begin{claim}
    For every $n$, $f_n$ is a bijection.
\end{claim}
\begin{proof}
    When $n=0$: $f_0:\{\varepsilon\}\to\{0\}$, $f_0(\varepsilon)=0$, is a
    bijection.

    Assume $f_{n-1}:\{0,1\}^{n-1}\to\{0,\dots,2^{n-1}-1\}$ is a bijection.

    For $f_n$,
    \[
        f_n(s_1\cdots s_n) = s_1\cdot 2^{n-1} + \sum_{i=2}^n s_i\cdot 2^{n-i}
        = s_1\cdot 2^{n-1} + f_{n-1}(s_2\cdots s_n).
    \]

    When $s_1=0$: $f_n(s_1\cdots s_n)\in\{0,\dots,2^{n-1}-1\}$.

    When $s_1=1$: $f_n(s_1\cdots s_n) \in \{2^{n-1},\dots,2^{n-1}-1+2^{n-1}\}
    = \{2^{n-1},\dots,2^n-1\}$.

    We know that $f_n(s_1\cdots s_n)\in\{0,\dots,2^n-1\}$.

    By induction, we know that $f_n$ is a bijection.
\end{proof}

\paragraph{Starting index of each length group.} Observe that, for
length $=n$, the starting index is equal to the number of bitstrings of
length $=0$ to length $=n-1$:
\[
    \text{starting index} = \sum_{k=0}^{n-1} |\{0,1\}^k| = \sum_{k=0}^{n-1} 2^k
    = \frac{2^n-1}{2-1} = 2^n-1.
\]

So, $f : \{0,1\}^* \to \N$ is equal to $(2^n-1) + $ (position of bitstring
in the group of the length):
\[
    f(s) = (2^n-1) + f_n(s), \qquad n = |s|.
\]

\begin{claim}
    $f$ is a bijection.
\end{claim}
\begin{proof}
    We know $f(s) \in \big[\,2^{|s|}-1,\ 2^{|s|+1}-2\,\big]$.

    \textbf{First, we prove $f$ is injective.}

    \emph{Case 1: $s\ne t$ and $|s|\ne|t|$.} Let $|s|<|t|$. Then
    \[
        f(s) \in \big[2^{|s|}-1,\ 2^{|s|+1}-2\big], \qquad
        f(t) \in \big[2^{|t|}-1,\ 2^{|t|+1}-2\big].
    \]
    Since $|s|<|t|$ are integers, $|s|+1\le|t|$, so
    \[
        2^{|s|+1}-2 \le 2^{|t|}-2 < 2^{|t|}-1.
    \]
    So $\big[2^{|s|}-1,\ 2^{|s|+1}-2\big] \cap \big[2^{|t|}-1,\ 2^{|t|+1}-2\big]
    = \varnothing$. Hence $f(t)\ne f(s)$.

    \emph{Case 2: $s\ne t$ and $|s|=|t|$.}
    \[
        f_{|s|}(s) \ne f_{|t|}(t) \implies f(s) \ne f(t).
    \]

    By Case 1 and 2, $f$ is injective.

    \textbf{Now we prove $f$ is surjective}, i.e.\ need to prove
    $\forall m\in\N,\ \exists\, s\in\{0,1\}^*,\ f(s)=m$.

    We know $f(s) \in \big[2^{|s|}-1,\ 2^{|s|+1}-2\big]$, so we need
    \[
        2^{|s|}-1 \le m \le 2^{|s|+1}-2
        \iff
        2^{|s|} \le m+1 \le 2^{|s|+1}-1
        \iff
        2^{|s|} \le m+1 < 2^{|s|+1}
        \iff
        |s| \le \log_2(m+1) < |s|+1.
    \]

    So let $n = \lfloor \log_2(m+1) \rfloor$. Then
    \[
        2^n \le m+1 < 2^{n+1}
        \implies
        0 \le m - 2^n + 1 < 2^n.
    \]

    So $m-2^n+1 \in [0,2^n) \subseteq [0,2^n-1]$, which is the domain of
    $f_n^{-1}$ (i.e., $m-2^n+1$ is in the range of $f_n$).

    $f_n$ is surjective, so $\exists\, s\in\{0,1\}^n$ such that
    $f_n(s) = m-2^n+1$. Then
    \[
        f(s) = (2^n-1) + f_n(s) = (2^n-1) + (m-2^n+1) = m.
    \]

    We proved that $\forall m\in\N$, $\exists\, s\in\{0,1\}^*$, $f(s)=m$.
    So, $f$ is surjective.

    Hence, $f$ is a bijection.
\end{proof}

\end{solution}
\end{solbox}

\begin{exercise}{9}
    Let $A$ be a set of $n$ elements. A function $f:A\to A$ \emph{fixes} $x$ if
    $f(x)=x$ ($x$ is a \emph{fixed point} of $f$). For $X\subseteq A$, $f$ fixes
    $X$ if it fixes every element of $X$, i.e.\ $f|_X = i_X$.

    Let $\mathrm{Fix}(X)$ be the set of bijections on $A$ that fix $X$, and let
    $F$ be the set of all bijections on $A$ that fix at least one element of $A$.
    Using Exercise 1.13, express $|F|$ in terms of the set sizes
    $|\mathrm{Fix}(X)|$ using all $X\subseteq A$.
\end{exercise}
\begin{solbox}
\begin{solution}

\end{solution}
\end{solbox}

\begin{exercise}{10}
    For each $k\in\{1,\dots,n\}$ and any $X\subseteq A$ with $|X|=k$, write $f_k$
    for the number of bijections on $A$ that fix $X$.
    \begin{enumerate}[label=(\alph*)]
        \item Why don't we include $X$ in the notation $f_k$?
        \item Give an expression for $f_k$.
        \item Express $|F|$ in terms of $f_1,\dots,f_n$, and then use your
              expression for $f_k$ to give an expression for $|F|$.
        \item Hence express the number of fixed-point-free bijections on $A$ in
              terms of $f_1,\dots,f_n$.
        \item For each of the following sets, what proportion of all bijections
              on the set are fixed-point-free?
              $\{0,1\}$; $\{\spadesuit,\clubsuit,\diamondsuit,\heartsuit\}$; the
              set of ten decimal digits; the 26-letter English alphabet.
        \item What is the largest set for which you can determine this
              proportion, and what is its value?
    \end{enumerate}
\end{exercise}
\begin{solbox}
\begin{solution}[(a)]

\end{solution}
\begin{solution}[(b)]

\end{solution}
\begin{solution}[(c)]

\end{solution}
\begin{solution}[(d)]

\end{solution}
\begin{solution}[(e)]

\end{solution}
\begin{solution}[(f)]

\end{solution}
\end{solbox}

\begin{exercise}{11}
    \begin{enumerate}[label=(\alph*)]
        \item Write down all surjections from $\{1,2,3,4\}$ to $\{a,b,c\}$.
              Compare your answer with Exercise 1.20.
        \item Define a surjection from
              \[
                  \{\text{surjections from } \{1,2,3,4\} \text{ to } \{a,b,c\}\}
              \]
              to
              \[
                  \{\text{partitions of } \{1,2,3,4\} \text{ into three parts}\}.
              \]
        \item What sizes can the preimages of members of its image be?
    \end{enumerate}
\end{exercise}
\begin{solbox}
\begin{solution}

\end{solution}
\end{solbox}

\begin{exercise}{12}
    Suppose $f:A\to B$ and $g:C\to D$ are bijections, $A\cap C=\emptyset$ and
    $B\cap D=\emptyset$. Show that $f\sqcup g$ is a bijection and express the
    inverse of $f\sqcup g$ in terms of the inverses of $f$ and $g$.
\end{exercise}
\begin{solbox}
\begin{solution}

\end{solution}
\end{solbox}

\begin{exercise}{13}
    If the function $f$ has an inverse function, which of its iterated
    compositions $f^{(k)}$ ($k\in\N$) have inverse functions?
\end{exercise}
\begin{solbox}
\begin{solution}

\end{solution}
\end{solbox}

\begin{exercise}{14}
    Determine all functions from $\{1,2,3\}$ to itself whose indefinitely
    iterated composition is a constant function (i.e.\ $f^{(k)}$ is constant for
    large enough $k$).

    Challenge: characterise those $f:\{1,\dots,n\}\to\{1,\dots,n\}$ such that, for
    large enough $k$, $f^{(k)}$ is constant.
\end{exercise}
\begin{solbox}
\begin{solution}

\end{solution}
\end{solbox}

\begin{exercise}{15}
    \textbf{Caesar slide cryptosystem.} Each letter of the message is slid
    rightwards along the alphabet (with wrap-around) by a fixed key $k$. E.g.\
    with $k=3$ (i.e.\ key letter \texttt{d}),
    \texttt{thefamilyofdashwoodhadlongbeen} encrypts to
    \texttt{wkhidplobrigdvkzrrgkdgorqjehhq}.

    Formally: message space $M$ = all strings of lower-case English letters; key
    space $K=\{\mathtt a,\dots,\mathtt z\}$; cypher space $C=M$. Identify letters
    with $\{0,\dots,25\}$ and let $+$/$-$ be letter addition/subtraction mod 26.
    \begin{itemize}
        \item Encryption $e:M\times K\to M$:
              $e(m_1\cdots m_n,k)=c_1\cdots c_n$ with $c_i=m_i+k$.
        \item Decryption $d:M\times K\to M$:
              $d(c_1\cdots c_n,k)=m_1\cdots m_n$ with $m_i=c_i-k$.
    \end{itemize}
    Using the definitions of \emph{embeddable}, \emph{equivalent} and
    \emph{idempotent} from Assignment 1, prove that the Caesar slide cryptosystem
    is idempotent.
\end{exercise}
\begin{solbox}
\begin{solution}

\end{solution}
\end{solbox}

\begin{exercise}{16}
    If $f$ is an injection and $f^{-1}$ is its inverse relation, what can you say
    about $f\circ f^{-1}$ and $f^{-1}\circ f$? What kinds of functions or
    relations are they, and what are their domains and codomains?
\end{exercise}
\begin{solbox}
\begin{solution}

\end{solution}
\end{solbox}

\begin{exercise}{17}
    If a binary relation is symmetric, what can you say about its inverse
    relation?
\end{exercise}

\begin{exercise}{18}
    Recall the Parent relation from \S2.13. Suggest a good name for its
    transitive closure. Then use that relation to define formally the relation
    that links two people that are related to each other (however distantly).
\end{exercise}
\begin{solbox}
\begin{solution}

\end{solution}
\end{solbox}

\begin{exercise}{19}
    Rephrase the Collatz Conjecture as a statement about the transitive closure
    of the Collatz function.
\end{exercise}
\begin{solbox}
\begin{solution}

\end{solution}
\end{solbox}

\begin{exercise}{20}
    Below are some binary relations on the set of all Python programs $P,Q$. For
    each relation, determine whether or not it is an equivalence relation, and
    give reasons.
    \begin{center}
    \begin{tabular}{ll}
        notation & the relation holds when\ldots\\\hline
        $P\simeq_1 Q$ & $P$ and $Q$ have the same number of characters\\
        $P\simeq_2 Q$ & $P$ and $Q$ compute the same function\\
        $P\simeq_3 Q$ & $P$ and $Q$ have the same set of variable names\\
        $P\simeq_4 Q$ & $P$ and $Q$ have at least one variable name in common\\
        $P\simeq_5 Q$ & $P$ and $Q$ have no variable names in common\\
        $P\simeq_6 Q$ & $P$ and $Q$ were written by the same set of programmers\\
        $P\simeq_7 Q$ & $P$ and $Q$ have at least one coauthor in common\\
        $P\simeq_8 Q$ & $P$ was completed before $Q$\\
        $P\simeq_9 Q$ & $P$ was completed on the same day as $Q$\\
        $P\simeq_{10} Q$ & $P$ was completed after $Q$
    \end{tabular}
    \end{center}
\end{exercise}
\begin{solbox}
\begin{solution}

\end{solution}
\end{solbox}

\begin{exercise}{21}
    What can you say about the transitive closure of an equivalence relation?
\end{exercise}

% ======================================================================
\begin{exercise}{22}
    The \emph{anagram} relation on English words: $(x_1\cdots x_m,\,y_1\cdots
    y_n)$ belongs to it iff there exists a bijection
    $f:\{1,\dots,m\}\to\{1,\dots,n\}$ such that $x_i=y_{f(i)}$ for all $i$.
    \begin{enumerate}[label=(\alph*)]
        \item Find the largest set of three-letter words you can that are all
              related to each other by anagram.
        \item For each $k=1,2,3,4$, find an English word of length $k$ that is
              not related to any other word by anagram.
        \item Prove that anagram is an equivalence relation.
        \item (Programming challenge) Using a standard open-access list of
              English words, determine the number of equivalence classes of
              anagram on that list.
    \end{enumerate}
\end{exercise}
\begin{solbox}
\begin{solution}[(a)]

\end{solution}
\begin{solution}[(b)]

\end{solution}
\begin{solution}[(c)]

\end{solution}
\end{solbox}

\begin{exercise}{23}
    How many non-reflexive binary relations are there on a set of size $n$?
\end{exercise}

\begin{exercise}{24}
    The rows of a table with $n$ columns represent the $n$-tuples of an $n$-ary
    relation on $A_1\times\cdots\times A_n$. Assuming the columns are all
    distinct, how many ternary relations can you construct by choosing columns
    from this table?
\end{exercise}
\begin{solbox}
\begin{solution}

\end{solution}
\end{solbox}

